\Def Контейнер — некоторая структура, позволяющая хранить
некоторые данные. Например, массив.

\Def Контейнер линейный, если на нем однозначно определено
отношение «следующий элемент».

\Def Список — последовательный набор узлов, предоставляющий
следующий интерфейс.

\begin{center}
   \centering
   \begin{tabular}{|c|c|c|}
    \hline
    Операция                        & Время & Примечание          \\ \hline
    Вставка в начало                & O(1)  &                     \\ \hline
    Удаление из начала              & O(1)  &                     \\ \hline
    Вставка в конец                 & O(1)  & Если двусвязный     \\ \hline
    Удаление из конца               & O(1)  & Если двусвязный     \\ \hline
    Вставка в произвольное место    & O(1)  & Если известно место \\ \hline
    Удаление из произвольного места & O(1)  & Если известно место \\ \hline
    Поиск                           & O(N)  &                     \\ \hline
    Обращение по индексу            & O(N)  &                     \\ \hline
    \end{tabular}
\end{center}

\Def Односвязный список -- последовательный набор из узлов с
данными, где каждый узел знает, где лежит следующий за ним.

\Def Двусвязный список -- последовательный набор из узлов с
данными, где каждый узел знает, где лежат следующий и предыдущий
узлы.

\Def Стек stack -- структура данных, которая хранит элементы и
предоставляет к ним доступ в рамках парадигмы LIFO (Last in, First
Out).
Можно реализовать на массиве и на односвязном списке.

\Def Очередь queue -- структура данных, которая хранит элементы и
предоставляет к ним доступ в рамках парадигмы FIFO (First in, First
Out).
Можно реализовать на двусвязном списке.

\Def Дек deque -- структура, представляющая из себя двустороннюю
очередь, то есть можно вставлять/удалять в начало/конец.
Можно реализовать на двусвязном списке.

\begin{center}
Стек с минимумом
\end{center}

При добавлении нового элемента в стек $S$ для получения стека $S'$
будем в голову $S'$ добавлять минимум как минимум из значения
элемента и значения минимума в голове $S'$.

\begin{center}
Очередь на двух стеках
\end{center}

Сделаем один стек на вход stack\_in и второй на выход stack\_out. В
первый будем добавлять, а из второго — извлекать.
Если при извлечении stack\_out пуст, то перекидываем все элементы
из stack\_in в stack\_out.

\begin{center}
Очередь c поддержкой операции получения минимума
\end{center}
Заведем очередь на двух стеках, поддерживающих минимум.
Запрос минимума сводится к минимуму из минимумов в двух стеках.

\pagebreak